% ============================================================
% 8.3 — kompakte Arbeitsfassung, Phase 02b, 10.09.2026.
% Forschungsfrage ohne "effizient" nach Autorenpräzisierung.
% Schluss aus 3–7 und 8.1/8.2; keine neuen Ergebnisse oder Quellen.
% Kein separates Kapitel 9 vorgesehen.
% ============================================================

\section{Fazit und Ausblick}
\label{sec:fazit-ausblick}

Die Forschungsfrage \emph{Wie können KI-Agenten mittels Microsoft Agent
Framework in den frühen Phasen des Software Development Lifecycle
eingesetzt werden?} beantwortet die Arbeit mit einem konkret
realisierten Integrationsansatz: KI-Agenten und modellgestützte Stufen
überführen Projektkommunikation in Anforderungen, Architekturinformationen
und Arbeitspakete. MAF verbindet diese Verarbeitung mit expliziten
Workflows, Werkzeugen und menschlichen Entscheidungspunkten. Quellenbindung
und ein laufübergreifender Core ermöglichen es, Inhalte im Zusammenhang
ihrer Herkunft, Beziehungen und Änderungen zu führen. Die Fortschreibung
wird durch fachliche Freigaben und technische Integritätsprüfungen
begrenzt. Die Implementierung des fachlichen Zielsystems liegt außerhalb
des untersuchten SDLC-Ausschnitts.

Der Beitrag umfasst das Artefakt, die Rekonstruktion seiner explorativ
begründeten Gestaltung und die Untersuchung ausgewählter Eigenschaften.
Die Befunde zeigen realisierte Möglichkeiten der Quellenverfolgung und
kontrollierten Fortschreibung ebenso wie unvollständige Abdeckung,
zusätzlichen Tokenaufwand und offene Herkunftspfade. Entscheidend für den
gewählten Ansatz ist die begründete Verteilung von semantischem
Entscheidungsspielraum, verbindlicher Ablaufsteuerung und menschlicher
Verantwortung. Ihre konkrete Umsetzung ist nachvollziehbar beschrieben;
ihre Leistungsfähigkeit bleibt auf den geprüften Umfang begrenzt.

Weitere Arbeit sollte zunächst die identifizierten Herkunftslücken und
offenen Kontrollfälle bearbeiten. Anschließend wäre ein konsolidierter
Systemstand mit wiederholten Ausführungen, weiteren Quellenfällen und
einer stärker unabhängigen menschlichen Inhaltsprüfung zu untersuchen.
Für den praktischen Einsatz ist besonders relevant, ob menschliche
Bearbeitung anhand der bereitgestellten Hinweise die Ergebnisqualität
verbessert und welcher Prüfaufwand dabei entsteht. Diese Untersuchung
würde die bislang getrennten Befunde zu maschineller Verarbeitung und
kontrollierter Übernahme um die tatsächliche Nutzung ergänzen.
